\chapter{Traumatologische Notfälle ---
\gAasRsN{RS~VII}}

Michael Withofner, Dr. med. univ.

\minitoc

\paragraph{Trauma}

\gDefinition{Als Trauma  bezeichnet man eine {\textperiodcentered}~Sch\"adigung,
{\textperiodcentered}~Verletzung, oder {\textperiodcentered}~Wunde, die
durch eine Gewalteinwirkung physikalischer oder chemischer Art
verursacht wurde.}

Die Traumata lassen sich grob ind folgende Klassen Einteilen:
% 
% \begin{itemize}
%     \item Sch\"adel-Hirn-Trauma (SHT)
%     \item Thorax trauma
%     \item Abdominal trauma / Bauchtrauma
%     \item Beckentrauma
%     \item Verletzungen der Wirbels\"aule
%     \item Polytrauma
%     \item Verletzungen der Extremit\"aten
%     \item Verbrennungen
% \end{itemize}



%%%%%%%%%%%%%%%%%%%%%%%%%%%%%%%%%%%%%%%%%%%%%%%%%%%%%%%%%%%
%%%%%%%%%%%%%%%%%%%%%%%%%%%%%%%%%%%%%%%%%%%%%%%%%%%%%%%%%%%
%%%%%%%%%%%%%%%%%%%%%%%%%%%%%%%%%%%%%%%%%%%%%%%%%%%%%%%%%%%

\section{Wunden}\label{topic:}\index{}

\subsection{Mechanische
Wunden}\label{topic:wunden-meschanische}\index{Wunden!mechanische}

\begin{figure}[H]
    \begin{center}
\includegraphics[width=5cm]{./images/stop.jpg}
\captionof{figure}{\gAuthNotesPicLegalNotOk{}{}Muster
\gSourcePicture{}}
\end{center}
\end{figure}




\begin{tabularx}{\linewidth}[H]{XXX}
\\\toprule
\multicolumn{3}{c}{{\gTabHeading{Diashow Muster}}}
\\\midrule
\includegraphics[width=\linewidth]{./images/stop.jpg}
\captionof{figure}{\gAuthNotesPicLegalNotOk{}{}Muster
\gSourcePicture{}} &
\includegraphics[width=\linewidth]{./images/stop.jpg}
\captionof{figure}{\gAuthNotesPicLegalNotOk{}{}Muster
\gSourcePicture{}} &
\includegraphics[width=\linewidth]{./images/stop.jpg}
\captionof{figure}{\gAuthNotesPicLegalNotOk{}{}Muster
\gSourcePicture{}} \\\midrule
\includegraphics[width=\linewidth]{./images/stop.jpg}
\captionof{figure}{\gAuthNotesPicLegalNotOk{}{}Muster
\gSourcePicture{}} &
\includegraphics[width=\linewidth]{./images/stop.jpg}
\captionof{figure}{\gAuthNotesPicLegalNotOk{}{}Muster
\gSourcePicture{}} &
\includegraphics[width=\linewidth]{./images/stop.jpg}
\captionof{figure}{\gAuthNotesPicLegalNotOk{}{}Muster
\gSourcePicture{}} \\\bottomrule
\end{tabularx}

\subsubsection{Stichwunden}\label{topic:stichwunden}\index{Stichwunden}

\begin{gWideParEnv}
\begin{tabularx}{\linewidth}[H]{XXX}
\\\toprule
\multicolumn{3}{c}{{\gTabHeading{Diashow}}}
\\\midrule
\includegraphics[width=\linewidth]{./images/pneumothorax_stichverletzung_diskret.jpg}
\captionof{figure}{\gAuthNotesPicLegalOk{}{}Stichverletzng:
unauffällig. \gSourcePicture{Hauer}} &
\includegraphics[width=\linewidth]{./images/stichverletzung_oberarm.jpg}
\captionof{figure}{\gAuthNotesPicLegalOk{}{}Stichverletzung
am Oberarm, Schwellung. \gSourcePicture{Hauer}} &
{leer}
\\\bottomrule
\end{tabularx}
\end{gWideParEnv}


\subsubsection{Schnittwunden}\label{topic:schnittwunden}\index{Schnittwunden}

\begin{figure}[H]
    \begin{center}
\includegraphics[width=5cm]{./images/pulsader_schnitt.jpg}
\captionof{figure}{\gAuthNotesPicLegalOk{}{}Schnittverletzung
Oberhalb der "`Pulsader"'.
\gSourcePicture{Hauer}}
\end{center}
\end{figure}

\subsubsection{Rißwunden}\label{topic:risswunden}\index{Rißwunden}

\subsubsection{Quetschwunden}\label{topic:quetschwunden}\index{Quetschwunden}

\subsubsection{Rißquetschwunden
(RQW)}\label{topic:rqw}\index{Rißquetschwunden}\index{RQW}

% \begin{gDiasEnv3}
% \includegraphics[width=\linewidth]{./images/wunde_rqw1.jpg}
% \captionof{figure}{\gAuthNotesPicLegalOk{}{}Rißquetschwunde
% vor der Versorgung im Spital.
% \gSourcePicture{Hauer}}
% &
% \includegraphics[width=\linewidth]{./images/wunde_rqw2.jpg}
% \captionof{figure}{\gAuthNotesPicLegalOk{}{}Rißquetschwunde
% nach der Versorgung im Spital.
% \gSourcePicture{Hauer}}
% &
% \end{gDiasEnv3}



\begin{figure}[H]
    \begin{center}
\includegraphics[width=5cm]{./images/wunde_rqw1.jpg}
\captionof{figure}{\gAuthNotesPicLegalOk{}{}Rißquetschwunde
vor der Versorgung im Spital.
\gSourcePicture{Hauer}}
\end{center}
\end{figure}

\begin{figure}[H]
    \begin{center}
\includegraphics[width=5cm]{./images/wunde_rqw2.jpg}
\captionof{figure}{\gAuthNotesPicLegalOk{}{}Rißquetschwunde
nach der Versorgung im Spital.
\gSourcePicture{Hauer}}
\end{center}
\end{figure}

\subsubsection{Schußwunden}\label{topic:schusswunden}\index{Schußwunden}




\subsection{Thermische
Wunden}\label{topic:wunden-thermische}\index{Wunden!thermische}

\subsubsection{Brandwunden}\label{topic:brandwunden}\index{Brandwunden}

\subsubsection{Erfrierungen}\label{topic:erfrierungen}\index{Erfrierungen}

\subsection{Chemische
Wunden}\label{topic:wunden-chemische}\index{Wunden!chemische}

\subsubsection{Verätzungen}\label{topic:veraetzungen}\index{Verätzungen}


%%%%%%%%%%%%%%%%%%%%%%%%%%%%%%%%%%%%%%%%%%%%%%%%%%%%%%%%%%%
%%%%%%%%%%%%%%%%%%%%%%%%%%%%%%%%%%%%%%%%%%%%%%%%%%%%%%%%%%%

\subsection{Stichverletzungen}

NIEMALS  HERAUSZIEHEN!!!!!!!!

steril abdecken + fixieren

verletzte Strukturen??, dient ev. als 
{\quotedblbase}St\"opsel{\textquotedblleft}??

Fremdk\"orper soll nicht tiefer rutschen



%%%%%%%%%%%%%%%%%%%%%%%%%%%%%%%%%%%%%%%%%%%%%%%%%%%%%%%%%%%
%%%%%%%%%%%%%%%%%%%%%%%%%%%%%%%%%%%%%%%%%%%%%%%%%%%%%%%%%%%
%%%%%%%%%%%%%%%%%%%%%%%%%%%%%%%%%%%%%%%%%%%%%%%%%%%%%%%%%%%

\section{Sch\"adel-Hirn-Trauma (SHT)}
Folge von Gewalteinwirkung auf den  Sch\"adel

\paragraph{Arten}

\begin{itemize}
\item Gedecktes SHT: ohne  Er\"offnung der harten Hirnhaut (Dura  mater)
\item Offenes SHT: mit  Er\"offnung der harten Hirnhaut (Dura  mater)
\end{itemize}

\paragraph{Begleitende Verletzungen}

\begin{itemize}
    \item Verletzung der Kopfschwarte  (Ri{\ss}quetschwunde)
    \item Frakturen von
    \begin{itemize}
        \item Sch\"adeldach
        \item Sch\"adelbasis (Blutung aus Ohren, Liquoraustritt  aus der Nase)
        \item Gesichtssch\"adel
        \item Unterkiefer
    \end{itemize}
\end{itemize}


\subsection{Gedecktes SHT}
 
\paragraph{Klassifikation}

\begin{itemize}
\item \gT{Gehirnersch\"utterung} -- SHT 1\textdegree --
Contusio cerebri
\item \gT{Gehirnprellung} -- 2\textdegree -- Contusio
cerebri
\item \gT{Gehirnquetschung} -- SHT 2\textdegree 
\end{itemize}




\subsubsection{Gehirnersch\"utterung (Commotio (cerebri) --
SHT 1\textdegree}
\index{Gehirnersch\"utterung}\index{Commotio (cerebri)}

\gMarried{Die Gehirnerschütterung ist die leichteste Form des
Schädelhirntraumas. Leit Symptom ist ein Trauma mit
anschließender kurzer Bewusstlosigkeit. Weiters kann es zu
flüchtigen neurologischen Ausfällen kommen, es kann auch eine
retrograde Amnesie (Gedächtnisverlust betreffenden Vorgänge
oder dem Unfall) bestehen. Eventuell lag der Patient über
Übelkeit erbrechen und Schwindel. Die Gehirnrschütterung
sieht keine bleibenden Schäden nach sich.}
{
\begin{itemize}
    \item  Bewu{\ss}tlosigkeit /Bewu{\ss}tseinstr\"ubung {\textless}  1h
    \item \"Ubelkeit, Erbrechen, Schwindel
    \item evt. fl\"uchtige neurologische Ausf\"alle
    \item retrograde  Amnesie  ({\quotedblbase}kann sich an den
    Unfallhergang nicht erinnern{\textquotedblleft} )
    \begin{itemize}
        \item im CT/MR [${\rightarrow}$R{}-\pageref{bkm:LexComputertomographie}]
        keine nachweisbaren Sch\"aden
        \item keine  bleibenden Sch\"aden
    \end{itemize}
\end{itemize}}






\subsubsection{Gehirnprellung (Contusio (cerebri) --
SHT 2\textdegree}

\index{Gehirnprellung}\index{Contusio (cerebri)}

\gMarried{Die Gehirnprellung ist eine schwerere Form des Schädel-Hirntraumas. 
Leitsymptom ist die \gEs{Bewusstlosigkeit größer
 als 1 h}, beziehungsweise die Bewusstseinseintrübung >24 h.
 Mittes weiterführender Diagnostik im Spital lassen sich die
 \gEs{Hirnschädigungen} nachweisen. Es kam zur
 Nervenverletzungen kommen, dies kann bis zu vegetativen Störungen führen.
 Bleibende Schäden sind möglich.}
 {
 \begin{itemize}
    \item Bewu{\ss}tlosigkeit  {\textgreater}  1h
    \item Bewu{\ss}tseinstr\"ubung {\textgreater}24h
    \item nachweisbare Gehirnsch\"adigung / bleibende  Ausf\"alle
    \item ev. Gehirnnervenabrisse 
    \item  vegetative St\"orungen
\end{itemize}}






\subsubsection{Gehirnquetschung -- SHT III{\textdegree}}

\index{Gehirnquetschung}

\gMarried{Die Gehirnquetschung ist eine sehr schwere Form des
Schädelhirntraumas. Es kommt zu einer \gEs{Bewusstlosigkeit
länger als 24 h} und zu teils sehr ausgeprägten
Begleitsymptomen wie zum Beispiel Hirndruckzeichen, Atem und
Kreislaufstörungen bis hin zum Atemstillstand durch ein Klemm
des Atemzentrums. 

Es kann weiteres zu intrakranielle
Blutungen kommen (epidural, subdural, subarachnoidal).
Außerdem kommt es zu verschiedenen neurologischen Ausfällen
und Symptomen wie zum Beispiel in einen oder auch zerebrale
Kampfanfälle. }
{
\begin{itemize}
    \item Bewu{\ss}tlosigkeit  {\textgreater}  24h
    \item Atmungs-/ Kreislaufst\"orungen
    \item Hirndruck symptomatik
    \item evt. Einklemmung des Atemzentrums ${\rightarrow}$
    zentraler Atemstillstand
    \item Behinderung der Hirndurchblutung
    \item Hirnblutungen
%     \begin{itemize}
%         \item epidural
%         \item subdural
%         \item subarachnoidal
%     \end{itemize}
    \item L\"ahmungen
    \item Zerebrale Kr\"ampfe
\end{itemize}}


Mit \gEs{bleibenden Schäden} und Langzeitfolgen ist
zu rechnen.

\gFazit{Das Leitsymtptom eines jeden SHT ist die
\gEs{Bewusstseinsstörung}.}


\subsection{Intrakranielle Blutungen}

\paragraph{\gT{Luzides Intervall} (lichtes Intervall)}

Die Symptome können eventuell erst Stunden nach der
Verletzung einsetzen. Dazwischen ist der Patient unauffällig.
Diese Zeitspanne wird Luzides Intervall genannt und ist
sehr gefährlich da man den Patienten leicht falsch
einschätzen kann.


\gAuthNotes{Zusammenfassen, keine Unterteilung in Sub-,
Epi- und Aubarachnoidalblutung, nur mehr "`Hirnblutung"'}


\subsubsection{Epidurale Blutung} \index{Epidurale Blutung}

\gDefinition{Blutung zwischen
der harten Hirnhaut und dem Schädelknochen.}

\gMarried{Bei der epiduralen Blutung kommt es zu einer Blutung zwischen
der harten Hirnhaut und dem Schädelknochen, bedingt durch
den Riss einer Hirnhaut-Arterie. Durch die Blutung steigt
der Hrindruck, es kommt dabei zu \gE{Hirndruckzeichen}.}
{\begin{itemize}
    \item Ri{\ss} einer Meningealarterie
    \item Einblutung zwischen Kalotte und Dura mater
    \item {\textquotedblleft}lichtes{\textquotedblright} (luzides )
    Intervall: einige Stunden
\end{itemize}}



\subsubsection{Subdurale  Blutung} \index{Subdurale Blutung}

\gDefinition{Blutung
unterhalb der harten Hirnhaut und oberhalt der weichen
Hirnhaut.}

\gMarried{Bei der subduralen Blutung kommend es zu einer ein Blutung
unterhalb der harten Hirnhaut. Der Verlauf ist oft
schleichend und kann sich über mehrere Tage erstrecken. Die
genaue Unterscheidung zwischen epiduraler und subduraler
Blutung ist präklinisch nicht möglich.}
{\begin{itemize}
    \item schleichender (!) Verlauf: oft \"uber Tage
    \item auch ohne  Trauma: Hypertoniker, Alkoholiker
\end{itemize}}



\subsubsection{Subarachnoidale Blutung -- SAB}
\index{Subarachnoidale Blutung} (\index{SAB}SAB)

\emph{Abkz.: \gT{SAB}}

% TODO test

\gDefinition{Bei der subarachnoidalen Blutung kommt es zu einer
Blutung unterhalb der Spinnengewebeshaut (Arachnoidea).} 

\gMarried{Dafür gibt es
unterschiedliche Ursachen, sie kann sowohl durch ein Trauma
ausgelöst werden als auch spontan entstehen. Bei der
\gE{spontanen} Subarachnoidalenblutung platzt typischerweise
ein bereits bestehendes kleines Aneurysma eines Hirngefäßes und
verursacht die Blutung. Dies kann aus heiterem Himmel
passieren.

\gFazit{Die subarachnoidalen Blutung kann durch ein Trauma
oder spontan entstehen. Sie kann ein traumatologischer oder
auch neurologischer Notfall sein.}

Leitsymptom ist der plötzliche, schwere,
\gEs{"donnerschlagartige" Kopfschmerz}, welcher mit
Hirndruckzeichen und einem oft raschen Verfall des
Patienten einhergeht. Der Patient ist akut \gE{vital
bedroht!}}
{
\begin{itemize}
    \item Trauma oder spontan
    \item z.B. geplatztes Aneurysma, Br\"uckenvenenabri{\ss} bei
    Decelerationstrauma
    \item Klassisch: pl\"otzlicher, schwerer,
    {\quotedblbase}donnerschlagartiger{\textquotedblleft} Kopfschmerz
    \item evt. Herd symptomatik, schneller
    Verfall,
\end{itemize}}






\subsection{Offenes SHT}

Hinweise:

\begin{itemize}
    \item gro{\ss}e RQW mit tast-/sichtbaren  Knochensplittern
    \item Hirnaustritt (schlechte Prognose)
    \item Blutungen aus Nase / Ohren
    \item Liquoraustritt, oft mit Blut vermischt
    ({\quotedblbase}Tupfertest{\textquotedblleft})
\end{itemize}


\subsection{Spezielle Diagnostik beim SHT}

Grundsätzlich ist beim Schäde-Hirn-Trauma die übliche
Diagnostik druchzuführen, die besonders beim SHT wichtigen
Punkte sind hier jedoch nochmal beschrieben.

Beurteilung von:

\paragraph{Neurologische Beurteilung}

Die Neurologische Beurteilung ist beim SHT natürlich
besonders wichtig und besonders genau durchzuführen.
Auffälligkeiten können Hinweise auf eine direkte
Hirnschädigung oder eine Hirnblutung sein. Ausserdem ist
der Verlauf sehr interessant.

\begin{description}
\item[B ewu{\ss}tseinslage] klar/getr\"ubt/bewu{\ss}tlos
\item[Motorik] 
\item[Pupillenweite] Hirndruckzeichen
\item[Pupillenreaktion] Hirndruckzeichen
\end{description}

\paragraph{Traumacheck} Ein Traumackeck ist obligat
und sollte keiner weiteren Erwähnung bedürfen.



\gAuthNotes{GABS: Die einzelnen Massnahmen in Bezug auf die
unterschiedlichen SHTs gehören definiert.}

\begin{gMassnahmenEnv}
    \paragraph{Spezielle Maßnahmen beim SHT --- Allgemein}
    
    \begin{itemize}
        \item Lagerung: leicht erhöhter Oberkörper, ca.
        30\,\textdegree , wenn keine Kontraindikation
        besteht.
        \item Wundversorgung
        \begin{itemize}
            \item Gedecktes SHT:
            \begin{itemize}
                \item Steriler Verband
            \end{itemize}
            \begin{itemize}
                \item steriler Verband locker auf offene
                Verletzungen
                \item \item Knochen / Hirnteile NICHT 
                zur\"uckstopfen /  entfernen!
            \end{itemize}
        \end{itemize}
    \end{itemize}
\end{gMassnahmenEnv}

\begin{gMassnahmenEnv}
    \paragraph{Spezielle Maßnahmen beim SHT ---
    Gehirnerschütterung}
    \begin{itemize}
        \item 
    \end{itemize}
\end{gMassnahmenEnv}

\begin{gMassnahmenEnv}
    \paragraph{Spezielle Maßnahmen beim SHT ---
    Gehirnquetschung}
    \begin{itemize}
        \item 
    \end{itemize}
\end{gMassnahmenEnv}

\begin{gMassnahmenEnv}
    \paragraph{Spezielle Maßnahmen beim SHT ---
    Gehirnerschütterung}
    \begin{itemize}
        \item 
    \end{itemize}
\end{gMassnahmenEnv}

\begin{gMassnahmenEnv}
    \paragraph{Spezielle Maßnahmen beim SHT ---
    Gehirnerschütterung}
    \begin{itemize}
        \item 
    \end{itemize}
\end{gMassnahmenEnv}



\begin{itemize}
\item (Gefahr: hohe Querschnittl\"ahmung)

\end{itemize}
\begin{itemize}
\item Sch\"adel/HWS/BWS bilden eine Linie

\end{itemize}
\begin{itemize}
\item Jedes SHT ist im Spital abzukl\"aren! (klinisch und  bildgebend)

\end{itemize}

\gCave{\gEs{Jedes SHT ist im Spital abzukl\"aren!} Es
besteht die Gefahr dass sich der Patient gegenw\"artig in der
\gEs{"`luciden Phase"'} 
befindet und sich sein Zustand erst in einigen Stunden, dann aber schlagartig verschlechtert. Im Zweifelsfall wird der
Patient dort zur  Beobachtung (24h lang) aufgenommen.}




%%%%%%%%%%%%%%%%%%%%%%%%%%%%%%%%%%%%%%%%%%%%%%%%%%%%%%%%%%%
%%%%%%%%%%%%%%%%%%%%%%%%%%%%%%%%%%%%%%%%%%%%%%%%%%%%%%%%%%%
%%%%%%%%%%%%%%%%%%%%%%%%%%%%%%%%%%%%%%%%%%%%%%%%%%%%%%%%%%%

\section{Thoraxtrauma}

\gDefinition{Verletzung des Brustkorbes und/oder der darin
enthaltenen Organe.}

\paragraph{Thoraxtrauma -- allgemeine  Symptome}

\begin{itemize}
    \item Dyspnoe
    \begin{itemize}
        \item Tachypnoe
        \item Zyanose
    \end{itemize}
    \item atemabh\"angige  Schmerzen,  Inspirationsstop
    \item ev. Hautemphysem (Luftansammlung in der Haut, typisches
    {\quotedblbase}Knistern{\textquotedblleft})
    \item Prellmarken (z.B. Lenks\"aule von PKW)
    \begin{itemize}
        \item ={\textgreater} gr\"undlicher Traumacheck!
    \end{itemize}
\end{itemize}



\paragraph{Thoraxtrauma -- spezifische  Symptome}

\begin{itemize}
    \item instabiler Thorax
    \begin{itemize}
        \item Sternumfraktur
        \item Serienrippenfraktur (zus. paradoxe Atmung --  Brustkorb senkt beim
        Einatmen)
    \end{itemize}
    \item Dyspnoe + Prellmarken +  Bluthusten
    \begin{itemize}
        \item Lungenprellung: Gasaustausch [F0EA?] , da  Lungenabschnitte
        gesch\"adigt
    \end{itemize}
    \item Brustschmerz, Schock
    \begin{itemize}
        \item Rhythmusst\"orungen ${\rightarrow}$ Verletzung des  Herzens
        \item schwerer Schock, schneller Verfall ${\rightarrow}$ Riss der  Aorta
        \item Blut im Schwall aus dem Mund ${\rightarrow}$
        Speiser\"ohrenverletzung
    \end{itemize}
\end{itemize}

\subsubsection{Pneumothorax}\label{topic:pneumothorax}\index{Pneumothorax}

\begin{center}
 [Warning: Image ignored] % Unhandled or unsupported graphics:
%\includegraphics[width=7.086cm,height=10.686cm]{AASdev20090228-img91.png}
\gTempPicMissing{Pneumothorax img91}
\end{center}

\begin{itemize}
    \item Lungen- (innerer) oder  Brustwandverletzung (\"au{\ss}erer)
    \begin{itemize}
        \item Luft gelangt in den Pleuraspalt
        \item Unterdruck aufgehoben
        \item Lunge kollabiert auf betroffener Seite
    \end{itemize}
    \item "`unkomplizierter"' Pneumothorax
    \begin{itemize}
        \item Luft kann ungehindert in Pleuraspalt und  wieder heraus
        \item {\quotedblbase}nur{\textquotedblleft} Ausfall des betroffenen
        Lungenfl\"ugels
    \end{itemize}
    \item Ventil-, \index{Spannungspneumothorax}Spannungspneumothorax
    \begin{itemize}
        \item Luft kann zwar hinein, aber nicht wieder  heraus
        \item Druck auf betroffener Seite steigt ${\rightarrow}$  Organe (Lunge,
        Herz) werden auf die andere Seite gedr\"uckt
        \begin{itemize}
            \item gesunde Seite wird verdr\"angt + Abknickung  gro{\ss}er
            Gef\"a{\ss}e!
        \end{itemize}
    \end{itemize}
\end{itemize}

\begin{center}
 [Warning: Image ignored] % Unhandled or unsupported graphics:
%\includegraphics[width=5.811cm,height=4.343cm]{AASdev20090228-img92.png}
\gTempPicMissing{Spannungspneu img92}
\end{center}
Therapie

Pneumothorax, unkompliziert:

\begin{gMassnahmenEnv}
    \paragraph{Spezielle Maßnahmen --- Pneumothorax}
    \begin{itemize}
        \item \gTxMassVitBes
        \begin{itemize}
            \item $O_2$: 10\,-\,15\,l
            \item Lagerung: Oberkörper hoch. Wenn stabile
            Seitenlage notwendig ist dann Lagerung auf die
            \gE{verletzte} Seite, damit die unverletzte
            Lunge  nicht unnötig belastet wird.
        \end{itemize}
        \item Bei offener Thoraxverletzung: \gEs{Nicht
        luftdicht verbinden!} Es kann sonst ein
        Spannungspneumothorax entstehen. Steril abdecken. 
    \end{itemize}
\end{gMassnahmenEnv}

\gZ{Der Notarzt kann mittels Drainage einen Pneumothorax
entlasten.}

\gAuthNotes{KOCH: 

Es fehlt:

Serienrippenfraktur erklären
Paradoxe Atmung
Gefahr von Blutverlust
Therapie

Ev. auch kurz etwas über STernumfraktur schreiben?}

\section{Bauchtrauma}
\index{Bauchtrauma}

\begin{itemize}
    \item offen (perforierend),
    \begin{itemize}
        \item  Er\"offnung der  Bauchh\"ohle
        \item Stich-, Pf\"ahlungsverletzung
        \item Schnittverletzung, {\quotedblbase}Platzbauch{\textquotedblleft}
    \end{itemize}
    \item geschlossen (stumpf)
    \begin{itemize}
        \item Schlag, Tritt
        \item Prellung durch Lenkstange bei PKW-Unfall
    \end{itemize}
\end{itemize}




\subsection{Offenes Bauchtrauma}

\begin{center}
 [Warning: Image ignored] % Unhandled or unsupported graphics:
%\includegraphics[width=6.543cm,height=4.839cm]{AASdev20090228-img93.png}
\gTempPicMissing{Offenes Bauchtrauma img93}
\end{center}
\begin{itemize}
\item ev. Austritt von Darmschlingen

\end{itemize}

\begin{gMassnahmenEnv}
    
    \paragraph{Spezielle Maßnahmen beim offenen Bauchtrauma}
    \begin{itemize}
        \item \gTxMassVitBes
        \begin{itemize}
            \item Lagerung: Bauchdecke entspannen , Knierolle
        \end{itemize}
        \item Darmteile nicht  zur\"uckstopfen ${\rightarrow}$ sonst
        Darmri{\ss}/Darmverschlu{\ss}
        \item steril abdecken, mit Ringerl\"osung feuchthalten
        \begin{itemize}
            \item Gefahr des Eindringens von Keimen  (Peritonitis/Sepsis!)
        \end{itemize}
    \end{itemize}
\end{gMassnahmenEnv}


\subsection{Geschlossenes Bauchtrauma}

\paragraph{Gefahren}

\begin{itemize}
    \item Gef\"a{\ss}verletzungen
    \begin{itemize}
        \item Aorta, Vena cava,...
        \item massive innere Blutungen ${\rightarrow}$ Volumenmangelschock
    \end{itemize}
    \item Zerrei{\ss}ung  von Organen (\gEs{Milz},
    \gEs{Leber})
    \begin{itemize}
        \item massive \gEs{innere Blutungen} ${\rightarrow}$
        Volumenmangelschock
    \end{itemize}
    \item Zerrei{\ss}ung  von Hohlorganen (Magen, Darm,Galleng\"ange )
    \begin{itemize}
        \item Austritt von Speisebrei, Magensaft, Stuhl, Gallensekret
        \begin{itemize}
            \item Bauchfellentz\"undung (Peritonitis) nach Std. bis  Tagen
            \item Sepsis (hohe Lethalit\"at), septischer Schock
        \end{itemize}
    \end{itemize}
\end{itemize}

\paragraph{Symptome}

\begin{itemize}
    \item starke  Bauchschmerzen
    \item brettharte  Bauchdecke  (Abwehrspannung, D\'efense )
    \item Leitsymptom: Schockzeichen
    \item beim Bewu{\ss}tlosen schwierig zu  diagnostizieren
    \begin{itemize}
        \item gr\"undlich nach \index{Prellmarken}\gEs{Prellmarken}
        suchen!
    \end{itemize}
\end{itemize}

\gFazit{Immer nach \gEs{Prellmarken} suchen. \gEs{Patienten
entkleiden!}}


\subsubsection{Milzruptur}\index{Milzruptur}

Die Milz ist ein teuflisches Organ. Sie besitzt rund um
das gut durchblutete Gewebe eine Bindegewebs-Kapsel.

Das
Schlimme daran: Das Gewebe kann reißen, die Kapsel bleibt
jedoch intakt. Es kommt zuerst nur zu einer geringfügigen
Einblutung, dem Patienten geht es zunächst gut. Nach
einiger Zeit reißt dann die Kapsel und es kommt zu einer
ungehemmten Blutung in den Bauchraum. \gEs{Der Patient
verfällt von einem Moment auf den anderen} und wird
schockig.

\paragraph{Primäre Milzruptur} Die Milz und die Kapsel
reißen gleichzeitig, es kommt zu einem unmittel merkbaren
Schockzustand.

\paragraph{Sekundäre Milzruptur} Die Kapsel reißt erst
einige Zeit nach dem Riß des Milzgewebes \gE{(Latenzzeit)}.
Es kommt zu einem \gEs{zeitlich versetztem} und oft
\gEs{unerwartetem Schockgeschehen}. Mitunter reißt die
Kapsel erst Stunden nach dem eigentlichen Unfall. \gEs{Daher
ist es wichtig jeden Patienten, bei dem auch nur der geringste
begründete Verdacht auf eine Milzverletzung besteht, einer
weiterführenden Untersuchung im Spital zuzuführen!}
\gZ{Mittels Ultraschall kann auf eine Milzblutung untersucht
werden.}

\begin{itemize}
    \item prim\"ar: sofort, Kapsel durchgerissen
    \item sekund\"ar: Sickerblutung, Kapsel anfangs noch intakt
    ${\rightarrow}$ Latenz bis Ri{\ss} der Kapsel
\end{itemize}

\gCave{Cave! {\quotedblbase}\index{Zweizeitige
Milzruptur}Zweizeitige Milzruptur{\textquotedblleft}}



\subsubsection{Weitere interne Verletzungen}

Der Bauchraum beherbergt viele, teilweise sehr gut
durchblutete Organe die bei einer Verletzung zu erheblichen
Blutverlusten führen können. Präklinisch äußern sich diese
Verletzungen durch den durch den Blutverlust hervorgerufenen
Schockzustand. Weitere hinweiszeichen sind Prellmarken oder
Schmerzen in der entsprechenden Gegend.

Grundsätzlich kann man ursächlich gegen diese Verletzung
wenig tun, im Vordergrund steht die Erkennung,
Schockbehandlung und der zügige Transport in ein
Traumazentrum.

\paragraph{Aortenruptur}\index{Aortenruptur}\label{topic:aortenruptur}Durch
einwirkende starke Kräfte kann es zu einem Riß in der Aorta
kommen. Besonders bei Stürzen aus großen Höhen kann dies
vorkommen. Da es sich bei der Aorta um das größte Gefäß des
Körpers handelt ist die Prognose äußerst schlecht.

\paragraph{Nierenruptur}\index{Nierenruptur}\label{topic:nierenruptur}
Die Niere als naturgemäß gut durchblutetes Organ kann im
Falle einer Verletzung auch massive Probleme machen.

\paragraph{Leberuptur} Die Leber ist auch gut durchblutet
und kann im Falle einer Verletzung große Probleme machen.


\subsection{Beckentrauma}

\begin{itemize}
    \item erhebliche Gewalteinwirkung  (Versch\"uttung, \"Uberrolltrauma)
    \item massive  Blutungen (bis 5 l in wenigen  Minuten
    \item Gefahr des Ausblutens (ev. ohne sichtbare  Blutungsquelle!)
    \item Mitverletzung von Organen im kleinen  Becken:
    \begin{itemize}
        \item Dickdarm ${\rightarrow}$ Blutung aus Anus
        \item Harnsystem ${\rightarrow}$ Blut im Harn / Harn in der
        Bauchh\"ohle
        \item (Ausnahme: Spontanruptur der gef\"ullten  Harnblase bei stumpfem
        Bauchtrauma)
        \item Hoden ${\rightarrow}$ H\"amatome am Skrotum und am  Damm
    \end{itemize}
\end{itemize}


\begin{gMassnahmenEnv}
    \paragraph{Spezielle Maßnahmen bei Beckentraumen}
    \begin{itemize}
        \item \gTxMassVitBes
        \begin{itemize}
            \item Lagerung: flach
            \item $O_2$: 10-15\,l
        \end{itemize}
        \item Schockbek\"ampfung
        \item Bei Beckenfraktur:
        \index{Beckengurt}\gEs{Beckengurt}, zur Not mit
        G\"urtel
    \end{itemize}
\end{gMassnahmenEnv}



%%%%%%%%%%%%%%%%%%%%%%%%%%%%%%%%%%%%%%%%%%%%%%%%%%%%%%%%%%%
%%%%%%%%%%%%%%%%%%%%%%%%%%%%%%%%%%%%%%%%%%%%%%%%%%%%%%%%%%%
%%%%%%%%%%%%%%%%%%%%%%%%%%%%%%%%%%%%%%%%%%%%%%%%%%%%%%%%%%%

\section{Wirbels\"aulenverletzungen}

\gAuthNotes{KOCH: Überarbeiten, evt etwas kürzen.}

\paragraph{Typische Unfallmechanismen}

\begin{itemize}
    \item Sturz aus gro{\ss}er H\"ohe
    \item Motorradunfall (Genickbruch)
    \item eingeklemmte Personen bei VU
    \item Schlag auf Kopf (Stauchung der HWS)
\end{itemize}

\paragraph{Gefahren}

\begin{itemize}
    \item Kompression/Durchtrennung des  R\"uckenmarks
    \item L\"ahmungen (Parese)
    \item Sensibilit\"atsst\"orungen (Par\"asthesie)
    \begin{itemize}
        \item inkomplett: eines von beidem
        \item komplett: beides
    \end{itemize}
\end{itemize}

traumatische Einblutung ins R\"uckenmark --- dadurch bedingtes \"Odem ={\textgreater} RM-Kompression

\subsubsection{Schleudertrauma, Peitschenschlag}

bei Auffahrunf\"allen

Zerrung des Halsmarks (bildgebend meist nicht nachweisbar)

Schmerzen im HWS-Bereich

Kopfschmerzen

Schwindel, Tinnitus

ev. neurologische Ausf\"alle (in schweren F\"allen)

Prophylaxe: Nackenst\"utze

Therapie: Immobilisation, Stifneck, (ev. O2 )

\subsubsection{Querschnittssydrom}

inkomplett: Motorik oder Sensorik  ausgefallen

komplett: beides ausgefallen

Symptome:

schlaffe L\"ahmung, Ausfall aller Reflexe  unterhalb des betroffenen
Segmentes

Blasen-, Darml\"ahmung ={\textgreater} Inkontinenz,  \"Uberlaufblase

Ausfall der Gef\"a{\ss}-und W\"armeregulation

\subsubsection{Spinaler Schock}

pl\"otzliche spinale Durchtrennung mit  schlagartigem Wegfall aller
zentralen  erregenden Impulse...

Gef\"a{\ss}regulation versagt

Gef\"a{\ss}tonus [F0EA?]

RR [F0EA?]  

Wirbels\"aulenverletzungen -- {\quotedblbase}hoher{\textquotedblleft} 
Querschnitt

L\"asionen auf H\"ohe C1/C2 (1. bzw. 2.  HWK)

Abtrennung des Atemzentrums in der Medulla  oblongata

Anpralltrauma (Motorrad unter Lastwagen)

{\quotedblbase}K\"opfler{\textquotedblleft} in seichtes Gew\"asser

Erh\"angen

zentrale Ateml\"ahmung (keine  Spontanatmung!)

L\"asionen oberhalb C4 (4. HWK)

Segmente, die Zwerchfell und Kehlkopf  innervieren mitbetroffen

motorische Ateml\"ahmung ={\textgreater} auf Dauer  beatmungspflichtig

\begin{gMassnahmenEnv}
    
    \gEs{Bei Verdacht} auf WS-Trauma erfolgt eine Strikte
    Immobilisierung (Vacuum-Matratze, Schaufeltrage, Stifneck)
    
    Bewegungen nur in der WS-Achse (keine  Beugung, Drehung,
    Seitw\"artsneigung)
    
    \gTxMassVitBes
    
    \begin{itemize}
        \item $O_2$ im Rahmen der Schockbehandlung
        \item Lagerung: flach.
        \item Wenn Schocklagerung notwendig:
        \begin{itemize}
            \item spezielle Technik: Trage schr\"ag: Kopf tief,
            \item Kontraindikation: SHT oder kardiale
            Komplikationen
        \end{itemize}
        \item Schockbekäpfung
        \item Traumacheck
        \item Neurocheck
    \end{itemize}
    
\end{gMassnahmenEnv}





%%%%%%%%%%%%%%%%%%%%%%%%%%%%%%%%%%%%%%%%%%%%%%%%%%%%%%%%%%%
%%%%%%%%%%%%%%%%%%%%%%%%%%%%%%%%%%%%%%%%%%%%%%%%%%%%%%%%%%%
%%%%%%%%%%%%%%%%%%%%%%%%%%%%%%%%%%%%%%%%%%%%%%%%%%%%%%%%%%%

\section{Polytrauma}

\gAuthNotes{Neufassung}

\marginlabel{Polytrauma}\gDefinition{Gleichzeitige
Verletzung \gEs{mehrerer K\"orperregionen} oder
Organsystemen, von denen \gEs{wenigstens eine} allein
\gEs{oder die Kombination mehrerer} \gEs{lebensgef\"ahrlich}
ist.}

Zum Beispiel:
\begin{itemize}
    \item Bauchtrauma + Oberschenkel; SHT + WS + Thoraxtrauma etc.
    \item gebrochener Finger + gebrochene Zehe ${\rightarrow}$ kein
    Polytrauma
\end{itemize}

\gFazit{Patient gilt als polytraumatisiert  bis zum Beweis 
des Gegenteils!} 

\begin{gMassnahmenEnv}
    \paragraph{Spezielle Maßnahmen bei Polytrauma}
    \begin{itemize}
        \item Schneller Traumacheck gleichzeitig  mit  Sichern der
        Vitalparameter!
        \item Bewu{\ss}tsein
        \item Atmung
        \item Kreislauf: Carotispuls, Pulse peripher  (orientierend!),
        Kapillardurchblutung
        \item gr\"undlicher Traumacheck erst sp\"ater
        \item stabilisierende Ma{\ss}nahmen
        \item Atemwege sichern (Intubation),  ven\"oser  Zugang (Notarzt!)
        \item wenn einigerma{\ss}en stabil, gr\"undlicher  Traumacheck,
        Verletzungsmuster (={\textgreater} zu  erwartende Komplikationen)
    \end{itemize}
\end{gMassnahmenEnv}


\subsubsection{Polytrauma beim Kind}

Der Schockindex ist beim Kind unbrauchbar, da die
Normalwerte andere sind!

Kinder bleiben lange stabil, verfallen dann aber schlagartig!

\subsubsection{Der Unfallort}

\begin{itemize}
    \item Absichern  der Unfallstelle,  m\"oglichst auff\"allig machen (Blaulicht,
 Warnblinkanlage, Scheinwerfer)
 \item bei mehreren Verletzten: Triage
 \item Nachalarmierung : andere  Rettungsmittel, Feuerwehr, Polizei
 \item Eingeklemmte im Fahrzeug stabilisieren,  nicht eigenm\"achtig bergen
 \item Crash -Rettung nur  bei  Fremdgef\"ahrdung, HWS stabilisieren!
\end{itemize}



%%%%%%%%%%%%%%%%%%%%%%%%%%%%%%%%%%%%%%%%%%%%%%%%%%%%%%%%%%%
%%%%%%%%%%%%%%%%%%%%%%%%%%%%%%%%%%%%%%%%%%%%%%%%%%%%%%%%%%%
%%%%%%%%%%%%%%%%%%%%%%%%%%%%%%%%%%%%%%%%%%%%%%%%%%%%%%%%%%%

\section{Extremit\"atenverletzungen}
Gelenke betreffend

Verstauchung (Distorsion )

Verrenkung (Luxation )

Knochenbruch (Fractur )

\subsection{Verstauchung}
\gDefinition{Zerrung / \"Uberdehnung des
Kapsel-/Bandapparates}

\paragraph{Symptome}
\begin{itemize}
    \item Schmerzen
    \item Schwellung
    \item Blutergu{\ss} (H\"amatom )
\end{itemize}



\begin{gMassnahmenEnv}
    \paragraph{Spezielle Massnahmen bei einer Verstauchung}
    \begin{itemize}
        \item ruhigstellen
        \item hochlagern
        \item k\"uhlen
        \item ad Unfallabteilung (mu{\ss} nicht sofort sein)
    \end{itemize}

    
    
\end{gMassnahmenEnv}

\subsection{Verrenkung}

\gDefinition{Teilweise oder vollst\"andige Verrenkung  eines
Gelenks}

\paragraph{Symptome}
\begin{itemize}
    \item Schmerzen
    \item Schwellung
    \item Blutergu{\ss}
    \item abnorme Gelenkstellung (federnd fixiert), bewegungsunf\"ahig
\end{itemize}



\begin{gMassnahmenEnv}
    \paragraph{Spezielle Maßnahmen bei einer Verrenkung}
    
    Stellung belassen
    
    ruhigstellen
    
    KEINE EINRENKVERSUCHE!!!
\end{gMassnahmenEnv}



%%%%%%%%%%%%%%%%%%%%%%%%%%%%%%%%%%%%%%%%%%%%%%%%%%%%%%%%%%%
%%%%%%%%%%%%%%%%%%%%%%%%%%%%%%%%%%%%%%%%%%%%%%%%%%%%%%%%%%%
%%%%%%%%%%%%%%%%%%%%%%%%%%%%%%%%%%%%%%%%%%%%%%%%%%%%%%%%%%%

\subsection{Knochenbr\"uche}

\paragraph{Arten:}

\begin{itemize}
    \item geschlossene Fraktur
    \item offene Fraktur
\end{itemize}

jeweils

\begin{itemize}
    \item komplett
    \item inkomplett
\end{itemize}

\paragraph{Gefahren}

\begin{itemize}
    \item starke Blutungen
    \item ev. bis Schock
    \item begleitende Nervenverletzungen
    \item Fettembolie (Fraktur langer R\"ohrenknochen)
    \item Infektionsgefahr (Osteomyelitis bei offenen  Frakturen)
\end{itemize}

Immer gr\"undlichen Traumacheck  durchf\"uhren, um keine 
Begleitverletzungen zu \"ubersehen!!!


\paragraph{Frakturzeichen} \index{Frakturzeichen}

Es gibt sichere und unsichere Frakturzeichen:

\noindent\begin{tabularx}{\linewidth}[H]{XX}
\toprule
\multicolumn{2}{c}{\gEs{Frakturzeichen}}
\\\midrule
\multicolumn{1}{c}{\gEs{Sichere}}
 &
\multicolumn{1}{c}{\gEs{Unsichere}}
\\\toprule
\begin{itemize}
    \item Abnorme Beweglichkeit
    \emph{("`Pseudogelenk"')}\footnote{"`Pseudogelenk"':
    \emph{"`Wie ein Gelenk"'.} Beim Pesudogelenk hat es den
    Anschein als würde en Gelenk bestehen (das nicht
    dorthon gehört).)}
    \item Fehlstellung
    \item Knochenreiben (Krepitus) \footnote{\so{NICHT}
    absichtlich ausprobieren!}
    \item Sichtbare freie Knochenteile (= offener Bruch)
\end{itemize}
 &
 \begin{itemize}
     \item Schmerz
     \item Schwellung
     \item H\"amatom
     \item gest\"orte Funktion / Bewegungsunf\"ahigkeit
 \end{itemize}
\\\bottomrule
\end{tabularx}
\captionof{table}{Frakturzeichen}

\begin{figure}[h]
    \begin{center}
\includegraphics[width=6cm]{./images/fraktur_fehlstellung_nativ.jpg}
\includegraphics[width=6cm]{./images/fraktur_fehlstellung_xray.jpg}
\caption{Fehlstellung -- In natura und in
Röntgendarstellung. \gSourcePicture{Hauer}\gAuthNotesPicLegalOk{}{}}
\end{center}
\end{figure}

\paragraph{Spezielle Maßnahmen}

Blutstillung

Wundversorgung

ev. NAW-Nachalarmierung (starke Schmerzen, gr\"o{\ss}erer Blutverlust,
Schock, relevanten Begleit\-verletzungen)

Ruhigstellung / Schienung (unter Zug)


%%%%%%%%%%%%%%%%%%%%%%%%%%%%%%%%%%%%%%%%%%%%%%%%%%%%%%%%%%%
%%%%%%%%%%%%%%%%%%%%%%%%%%%%%%%%%%%%%%%%%%%%%%%%%%%%%%%%%%%
%%%%%%%%%%%%%%%%%%%%%%%%%%%%%%%%%%%%%%%%%%%%%%%%%%%%%%%%%%%
%%%%%%%%%%%%%%%%%%%%%%%%%%%%%%%%%%%%%%%%%%%%%%%%%%%%%%%%%%%

\section{Verbrennungen}

\gDefinition{\gAuthNotes{Definition fehlt}}

Die Gr\"o{\ss}e des gesch\"adigten Areals wird in Prozent der
ver\-brannten K\"orperoberfl\"ache angegeben.

\paragraph{Neunerregel}\label{topic:neunerregel}\index{Neunerregel}

Handfl\"ache (des Patienten) = 1~\% der K\"orperoberfl\"ache (KOF)

\gMarried{
\begin{center}
%\includegraphics[width=6.591cm,height=7.03cm]
\gAuthNotes{./images/AASdev20090228-img94.png}

\end{center}
}{
\begin{center}
\begin{tabulary}{\linewidth}[H]{LCC}
\toprule
 & \% &
\\\midrule
Kopf:
 & 9  &
\\
Arm, je:  &  9  &
\\
Thorax vorne: & 18  &
\\
Thorax hinten: & 18 &
\\
Bein, je: & 18 &
\\
Genital:  & 1 &
\\\bottomrule
\end{tabulary}
\end{center}
}

\paragraph{Gradeinteilung}

In Abhängigkeit von der Tiefe der Vebrennung wird das
Ausma{\ss} der Sch\"adigung durch Verbrennung wird in 3 Grade eingeteilt:

\begin{center}
\begin{tabulary}{\linewidth}[H]{CLC}
\toprule
 Grad
 &
 &
\\\midrule
1\,\textdegree
 &
R\"otung, Schmerz, Schwellung
 &
\\
2\,\textdegree
 &
Blasenbildung, Schmerz
 &
\\
3\,\textdegree
 &
Verkohlung, \gT{Verschorfung}, Schmerzen
abhängig von der teife der Schädigung
(Zerstörung von Nerven), Schrumpf\-ung des Gewebes durch
Hitze\-ein\-wirkung & \\\bottomrule
\end{tabulary}
\end{center}

\gCave{Ab \gEs{5\%} verbrannter (mind. 2. Grad) KOF bei
Kindern oder \gEs{10\%} bei Erwachsenen besteht Schockgefahr
(Fl\"ussigkeitsverlust)}


\begin{gMassnahmenEnv}
    \paragraph{Spezielle Ma{\ss}nahmen bei leichter
    Vebrennung}
    
    \begin{itemize}
        \item Eigenschutz!
        \item Kleiderbrand l\"oschen
        \item Kleidung rasch entfernen (sofern nicht  eingeschmolzen!)
        \item mind. 10 min. mit m\"a{\ss}ig kaltem Wasser d
       ie betroffenen Körperregionen spülen  --
       \gEs{NICHT unterk\"uhlen!}
        \item keimfreier, nicht verklebender Verband, locker anlegen
        \item wenn vorhanden: Water-Jel
        \footnote{Streitfrage: Water-Jel: Umstritten}
        \item Schockbek\"ampfung
        \item NAW wenn n\"otig
    \end{itemize}
    \gCave{CAVE Unterkühlung bei zu langer oder zu kalter
    Spülung!}
\end{gMassnahmenEnv}

\begin{gMassnahmenEnv}
    \paragraph{Spezielle Ma{\ss}nahmen bei großflächiger
    Vebrennung}
    \begin{itemize}
        \item \gTxMassVitBes
        \begin{itemize}
            \item $O_2$: 10\gLMin
        \end{itemize}
        \item Kleidung rasch entfernen (sofern nicht  eingeschmolzen!)
        \item mind. 10 min. mit m\"a{\ss}ig kaltem Wasser d
       ie betroffenen Körperregionen spülen  --
       \gEs{NICHT unterk\"uhlen!}
        \item keimfreier, nicht verklebender Verband, locker anlegen
        \item wenn vorhanden: Water-Jel
        \footnote{Streitfrage: Water-Jel: Umstritten}
        \item Schockbek\"ampfung
    \end{itemize}
    \gCave{CAVE Unterkühlung bei zu langer oder zu kalter
    Spülung!}
\end{gMassnahmenEnv}



\paragraph{Was man besser \gE{nicht} tun sollte \ldots}

\begin{itemize}
    \item NEIN: Salbe, Puder, etc. auf frische  Brandwunde
    \item NEIN: Blasen \"offnen
    \item NEIN: festklebende Kleidung gewaltsam  entfernen
    \item NEIN: aluminiumbeschichteten Verband  verwenden (Gefahr der
    Hitzereflexion bei  nicht ausreichend gek\"uhlten  Verbrennungen!)
\end{itemize}



\subsection{Inhalationstrauma}

\paragraph{Leitsymptome}

\begin{itemize}
    \item angebrannte Gesichtshaare
    \item geschwollene, ru{\ss}ige Lippen
    \item Ru{\ss}partikel im Auswurf
    \item Heiserkeit
    \item Stridor
\end{itemize}

${\rightarrow}$ Rauchgasvergiftung m\"oglich!

\begin{gMassnahmenEnv}
    \paragraph{Spezielle Maßnahmen beim Inhalationstrauma}~
    \gAuthNotes{fehlt!}
\end{gMassnahmenEnv}



%%%%%%%%%%%%%%%%%%%%%%%%%%%%%%%%%%%%%%%%%%%%%%%%%%%%%%%%%%%
%%%%%%%%%%%%%%%%%%%%%%%%%%%%%%%%%%%%%%%%%%%%%%%%%%%%%%%%%%%
%%%%%%%%%%%%%%%%%%%%%%%%%%%%%%%%%%%%%%%%%%%%%%%%%%%%%%%%%%%

\section{Erfrierungen}
\gDefinition{...sind Gewebsnekrosen, die durch lokale 
Durchblutungsst\"orungen aufgrund von  K\"alteeinwirkung (Eis, Wind, N\"asse,  K\"alte)
entstanden sind.}

Besonders geff\"ahrdet sind k\"orperferne K\"orperteile wie Zehen,
Finger, Ohren oder Nase.

Einteilung erfolgt in 4 Stadien, diese haben am Notfallort keine
entscheidende Bedeutung f\"ur die zu setzenden Ma{\ss}nahmen. Den
Schweregrad einer Erfrierung kann man erst nach Tagen voll
einsch\"atzen. Deshalb ist es wichtig, die Warnsymptome ernst zu
nehmen!

\paragraph{Stadien}

\gAuthNotes{fehlt!}

\paragraph{Symptome}

\begin{itemize}
    \item zuerst
    \begin{itemize}
        \item Gef\"uhllosigkeit und Bl\"asse
        \item starke Schmerzen und bl\"auliche Verf\"arbung
    \end{itemize}
    \item einige Zeit sp\"ater
    \begin{itemize}
        \item  marmorierte Verf\"arbung der Haut mit  Blasenbildung
        \item Taubheitsgef\"uhl, keine Empfindung bei  Ber\"uhrung
        \item Bewegungsunf\"ahigkeit des betroffenen  K\"orperteils
        \item starke Schmerzen
    \end{itemize}
\end{itemize}


\begin{gMassnahmenEnv}
    \paragraph{Spezielle Ma{\ss}nahmen bei Erfrierungen}~
\begin{itemize}
\item beengende Kleidungsst\"ucke lockern (Einschn\"urung erh\"oht
Risiko einer Erfrierung)
\item Keimfreier, gepolsterter Verband
\item Warme, gezuckerte Getr\"anke -- KEIN Alkohol
\item K\"orperstamm w\"armen, nie das erfrorene K\"orperteil abreiben!!
\end{itemize}
\end{gMassnahmenEnv}




%%%%%%%%%%%%%%%%%%%%%%%%%%%%%%%%%%%%%%%%%%%%%%%%%%%%%%%%%%%
%%%%%%%%%%%%%%%%%%%%%%%%%%%%%%%%%%%%%%%%%%%%%%%%%%%%%%%%%%%
%%%%%%%%%%%%%%%%%%%%%%%%%%%%%%%%%%%%%%%%%%%%%%%%%%%%%%%%%%%

\section{Unf\"alle durch Strom- und Blitzschlag}

\marginpar{\includegraphics[width=\linewidth]{./images/electricity_japan.jpg}
\small\emph{Foto: Fina}}

\gT{Niederspannung}: bis 1000\,Volt (Steckdose:  230 Volt)

\gT{Hochspannung}: \"uber 1000\,Volt, Stromst\"arke
{\textgreater}  3 Ampere (\"Uberlandleitungen) 

Die Verletzungsmuster variieren je nach  Spannung
(Volt) und Stromst\"arke (Ampere).

\begin{itemize}
    \item Reizeffekte (Muskel-/Nervenerregung):
    \begin{itemize}
        \item Verkrampfung ({\quotedblbase}kann nicht
        loslassen{\textquotedblleft})
        \item Sensibilit\"atsst\"orungen
        \item ZNS ${\rightarrow}$ Bewu{\ss}tseinsst\"orung
        \item Herzrhythmusst\"orungen (Kammerflimmern, ev.  als Sp\"atfolge!)
    \end{itemize}
    \item W\"armeentstehung
    \begin{itemize}
        \item Strommarken (Ein-/Austrittsverbrennung)
        \item Lichtbogen: oberfl\"achliche und tiefe  Verbrennungen
    \end{itemize}
\end{itemize}

\begin{itemize}
    \item Gleichstrom
    \item Wechselstrom ${\rightarrow}$ eher  Herzrhythmusst\"orungen
    \item Patient  kann Spannungsquelle ev. nicht  loslassen (Krampf der
    Handmuskulatur)
    \item Auch Defi stellt Strom-Gefahrenquelle  dar!!!
    \item Blitzunfall
    \begin{itemize}
        \item mehrere Millionen (!) Volt
        \item bis 100.000 Ampere
    \end{itemize}
\end{itemize}

\paragraph{Symptome}

\begin{itemize}
    \item Bewu{\ss}tlosigkeit
    \item L\"ahmungen (vor\"ubergehend)
    \item Blitzfiguren auf der Haut
    \item Herzrhythmusst\"orungen m\"oglich
\end{itemize}

\gFazit{Das eigentliche Ausma{\ss} der Verletzung ist oft
{\quotedblbase}unsichtbar{\textquotedblleft}, das der Strom im
K\"orperinneren -- also {\quotedblbase}unterirdisch{\textquotedblleft}
- viel Schaden anrichtet. }

\begin{gMassnahmenEnv}
    \paragraph{Spezielle Ma{\ss}nahmen}
    \begin{itemize}
        \item Eigenschutz!
        \begin{itemize}
            \item (Stromkreis  vorher  unterbrechen, FI bzw. bei Starkstrom durch
            Fachmann)
        \end{itemize}
        \item Genaue  Patientenuntersuchung ${\rightarrow}$ Eintritts{}-  und
        Austrittsmarke
        \item auf ev. Herzrhythmusst\"orungen vorbereitet  sein ${\rightarrow}$
        bis 24h Latenz
        \item Vitalparameter ${\rightarrow}$ entsprechende  symptomatische
        Versorgung
        \item NAW ${\rightarrow}$ Hospitalisierung
        \item Stromverbrennungen wie Verbrennungen  versorgen
    \end{itemize}
\end{gMassnahmenEnv}



